\documentclass{article}
\usepackage[T1]{fontenc}
\usepackage{lmodern}
\usepackage{graphicx}
\usepackage{amsmath,amssymb}
\usepackage[a4paper,margin=2cm]{geometry}
\usepackage[absolute,overlay]{textpos}
\setlength{\TPHorizModule}{1mm}
\setlength{\TPVertModule}{1mm}
\usepackage[indonesian]{babel}
\usepackage{hyperref}
\setlength{\emergencystretch}{2em}
% unit: Halaman 2

\begin{document}

% === Halaman 2 (llm) ===

\section*{Pendahuluan}

\begin{itemize}
    \item Kriptografi klasik (\textit{classical cryptography}) merupakan kriptografi yang sudah tua, sudah ada sejak ribuan tahun yang lalu hingga ditemukan komputer digital
    \item \textit{Cipher} pada kriptografi klasik, dinamakan \textit{cipher} klasik, (\textit{classical cipher}) hanya memproses pesan dari huruf-huruf alfabet saja
    \item Menggunakan alat tulis pena dan kertas saja, belum ada komputer
    \item Termasuk ke dalam jenis kriptografi kunci-simetri
\end{itemize}

\begin{itemize}
    \item Tiga alasan mempelajari kriptografi klasik:
    \begin{enumerate}
        \item Memahami konsep dasar kriptografi.
        \item Sebagai dasar algoritma kriptografi modern.
        \item Untuk memahami kelemahan sistem \textit{cipher}.
    \end{enumerate}
\end{itemize}

\end{document}
